\documentclass[conference]{IEEEtran}
\usepackage{shared/template_tgir_article}
\usepackage{booktabs}
\usepackage{longtable}
\usepackage{tabularx}
\usepackage{array}
\usepackage{enumitem}
\usepackage{xcolor}
\usepackage{listings}

\lstset{
    basicstyle=\ttfamily\small,
    breaklines=true,
    columns=fullflexible,
    frame=single
}

\newcommand{\product}{TGIR QuantLib Tools}
\newcommand{\code}[1]{\texttt{#1}}
\newcommand{\docdate}{2026-04-06}

% Visual design is controlled exclusively by template_tgir_article.sty and
% IEEEtran. Keep this file limited to content-support packages and commands.


\usepackage{amsmath,amssymb,mathtools}
\usepackage{microtype}
\usepackage{textcomp}
\usepackage{bookmark}
\usepackage{xurl}
\newcommand{\passthrough}[1]{#1}
\newcounter{none} % for unnumbered tables
\setlength{\emergencystretch}{3em} % prevent overfull lines
\providecommand{\tightlist}{}
\urlstyle{same}
\hypersetup{
  pdftitle={Callable Bermudan Cross-Currency Interest-Rate Swap Pricing Specification},
  pdfauthor={Tim Glauner, TG Investments and Research, LLC},
  pdfcreator={LaTeX via pandoc}}

\title{Callable Bermudan Cross-Currency Interest-Rate Swap Pricing
Specification}
\author{\TGIRArticleAuthorBlock}
\TGIRSetArticlePublicationDate{August 2026}

\begin{document}
\maketitle

{
\setcounter{tocdepth}{3}
}
\begin{center}
\begin{tabularx}{\columnwidth}{>{\bfseries}p{0.28\columnwidth}X}
\toprule
Status & Proposed specification---no implementation \\
Audience & Quantitative developers, model validators, and trading-system integrators \\
Last updated & 2026-08-14 \\
\bottomrule
\end{tabularx}
\end{center}

\section{Executive decision}\label{1-executive-decision}

Price the product with a three-factor Monte Carlo model and a two-pass
Longstaff-Schwartz least-squares Monte Carlo (LSM) exercise algorithm:

\begin{enumerate}
\def\labelenumi{\arabic{enumi}.}
\tightlist
\item
  one Hull-White one-factor short-rate model for the domestic currency;
\item
  one Hull-White one-factor short-rate model for the foreign currency;
  and
\item
  one lognormal Black-Scholes diffusion for the FX rate, correlated with
  both rate factors.
\end{enumerate}

This is a Monte Carlo model, not the short-rate lattice currently used
by the repository\textquotesingle s single-currency Bermudan swaptions.
That lattice uses Hull-White 1F by default and also supports a G2++
comparison. A multidimensional cross-currency lattice is possible in
principle, but it is not the recommended integration: it is harder to
build, calibrate, and converge, and QuantLib\textquotesingle s
\passthrough{\lstinline!TreeSwaptionEngine!} accepts one short-rate
model and one \passthrough{\lstinline!Swaption!} rather than a callable
cross-currency cash-flow portfolio.

The existing QuantLib integration remains the benchmark and supplies
conventions, curves, schedules, Hull-White calibration patterns, and
regression tests. The new pricer should be a separate engine so the
current \passthrough{\lstinline!ql.TreeSwaptionEngine!} results remain
unchanged.

\section{Product definition and sign
convention}\label{2-product-definition-and-sign-convention}

\subsection{Currencies and FX
quotation}\label{21-currencies-and-fx-quotation}

\begin{itemize}
\tightlist
\item
  \passthrough{\lstinline!DOM!} is the reporting, collateral, and NPV
  currency in the first implementation.
\item
  \passthrough{\lstinline!FOR!} is the other leg currency.
\item
  \passthrough{\lstinline!S(t)!} is quoted as units of
  \passthrough{\lstinline!DOM!} per one unit of
  \passthrough{\lstinline!FOR!}. A \passthrough{\lstinline!FOR!} cash
  flow \passthrough{\lstinline!CF\_FOR!} paid at time
  \passthrough{\lstinline!t!} is worth
  \passthrough{\lstinline!S(t) * CF\_FOR!} in
  \passthrough{\lstinline!DOM!} on that path.
\item
  All signed cash flows are from the option holder\textquotesingle s
  perspective: receipts are positive and payments are negative.
\end{itemize}

\subsection{Underlying swap}\label{22-underlying-swap}

The underlying contains two independently scheduled legs plus optional
notional exchanges:

\begin{itemize}
\tightlist
\item
  a domestic fixed, term-floating, or compounded-overnight leg;
\item
  a foreign fixed, term-floating, or compounded-overnight leg;
\item
  optional initial and final notional exchanges;
\item
  optional mark-to-market FX-reset notionals, explicitly out of scope
  for the first implementation unless selected as a separate product
  variant; and
\item
  spreads, gearing, stubs, payment lags, calendars, day counts, and
  business-day conventions on each leg.
\end{itemize}

The first supported trade should be a constant-notional cross-currency
swap with optional initial/final exchanges. Do not silently interpret a
mark-to-market cross-currency swap as constant-notional.

\subsection{Callable versus enterable Bermudan
rights}\label{23-callable-versus-enterable-bermudan-rights}

The request must carry \passthrough{\lstinline!exercise\_style!}:

\begin{itemize}
\tightlist
\item
  \passthrough{\lstinline!CANCEL\_REMAINING\_SWAP!} - the holder may
  terminate the remaining swap on specified dates. This is the initial
  target.
\item
  \passthrough{\lstinline!ENTER\_REMAINING\_SWAP!} - the holder may
  enter the remaining swap on specified dates, economically a
  cross-currency Bermudan swaption. This is a later compatible variant.
\end{itemize}

The request must also identify
\passthrough{\lstinline!exercise\_holder!} and the call settlement
amount from that holder\textquotesingle s perspective. A bare word such
as "callable" is insufficient because payer/receiver and issuer/holder
usages differ between systems.

For a cancellation right, exercising removes all contractually
cancellable future cash flows and replaces them with the defined
settlement cash flow. If the settlement is zero, a holder cancels when
continuing has sufficiently negative value. For an entry right,
exercising creates the remaining swap and any entry settlement.

\subsection{Exercise event
ordering}\label{24-exercise-event-ordering}

Each exercise opportunity has four dates:

\begin{itemize}
\tightlist
\item
  notice date;
\item
  decision date;
\item
  effective termination or entry date; and
\item
  settlement date.
\end{itemize}

The trade definition must state whether coupons and notional exchanges
on the effective date survive exercise. The default is: determine
already-fixed coupons first, pay cash flows scheduled on or before the
effective date, then remove later cancellable cash flows. No same-day
ordering may be inferred from date equality alone.

For version 1, \passthrough{\lstinline!notice\_date!} must equal
\passthrough{\lstinline!decision\_date!}; the timestamp at which notice
becomes irrevocable is the information time for the exercise decision.
Reject a trade with an earlier notice deadline until the state observed
at that deadline and the delayed exercise mechanics are specified.

\subsection{Decision-date cash-flow
partition}\label{25-decision-date-cash-flow-partition}

For a cancellation opportunity decided at
\passthrough{\lstinline!t\_i!}, partition domestic-equivalent
decision-date value into:

\begin{itemize}
\tightlist
\item
  \passthrough{\lstinline!A\_i!} - cash flows that survive either
  decision, including coupons/notional flows through the effective date
  under the trade\textquotesingle s event ordering;
\item
  \passthrough{\lstinline!J\_i!} - exercise-only settlement cash flows,
  converted and discounted from their actual settlement dates to
  \passthrough{\lstinline!t\_i!}; and
\item
  \passthrough{\lstinline!Y\_i!} - cancellable cash flows and later
  optionality that remain only if the holder continues.
\end{itemize}

The exercise comparison excludes the common component:

\begin{align*}
  \text{exercise if}\quad
  J_i &> \mathbb{E}^{\mathbb{Q}_d}\!\left[Y_i\mid\mathcal{F}_{t_i}\right]
        + \varepsilon_{\mathrm{ex}}, \\
  V_i &= A_i + \max\!\left(J_i,C_i\right), \\
  C_i &= \mathbb{E}^{\mathbb{Q}_d}\!\left[Y_i\mid\mathcal{F}_{t_i}\right].
\end{align*}

This partition applies even when the effective or settlement date is
later than the decision date. For a zero-fee cancellation,
\passthrough{\lstinline!J\_i = 0!}; surviving interim flows remain in
\passthrough{\lstinline!A\_i!} and do not disappear. For
\passthrough{\lstinline!ENTER\_REMAINING\_SWAP!}, define the exercise
value as the state-dependent PV of the entered swap plus settlement, and
compare it with the continuation value of the unexercised option.

\section{Scope}\label{3-scope}

\subsection{In scope}\label{31-in-scope}

\begin{itemize}
\tightlist
\item
  two currencies and one FX pair;
\item
  domestic-collateral valuation in the domestic currency;
\item
  Hull-White 1F dynamics for both short rates;
\item
  Black-Scholes lognormal FX dynamics with deterministic piecewise FX
  volatility;
\item
  a full constant 3x3 instantaneous correlation matrix, with an
  extension point for piecewise-constant correlation;
\item
  deterministic notionals in each currency;
\item
  discrete Bermudan cancellation dates;
\item
  LSM early exercise;
\item
  NPV, Monte Carlo error, cash-flow diagnostics, calibration
  diagnostics, exercise statistics, and bump-and-revalue risks; and
\item
  reproducible results from recorded market, trade, model, and
  random-seed inputs.
\end{itemize}

\subsection{Out of scope for the first
implementation}\label{32-out-of-scope-for-the-first-implementation}

\begin{itemize}
\tightlist
\item
  stochastic basis, stochastic volatility, local volatility, jumps,
  wrong-way credit risk, XVA, and funding valuation adjustments;
\item
  collateral switching or optional collateral currency;
\item
  mark-to-market/resettable notionals;
\item
  multiple exercise holders, partial calls, volume choice, or exercise
  into an amended coupon;
\item
  historical fixing ingestion beyond values supplied in the request;
\item
  automatic market-data sourcing; and
\item
  replacing the existing single-currency Bermudan tree pricer.
\end{itemize}

\section{Required input contract}\label{4-required-input-contract}

\subsection{Trade fields}\label{41-trade-fields}

At minimum:

\begin{center}
\begin{tabularx}{\columnwidth}{>{\bfseries}p{0.28\columnwidth}X}
\toprule
Group & Required fields \\
\midrule
Identity & \code{trade\_id}, \code{version}, \code{book}, and
  \code{as\_of\_date} \\
Product & \code{exercise\_style}, \code{exercise\_holder},
  \code{reporting\_currency}, and \code{collateral\_currency} \\
Domestic leg & Currency, pay/receive, notional, coupon type, rate or
  spread, index, schedule, day count, and payment lag \\
Foreign leg & The same fields, specified independently \\
FX & Pair, quote direction, spot, initial exchange rule, and final
  exchange rule \\
Exercise & Notice, decision, effective and settlement dates; settlement
  amount and currency; and same-day event order \\
Fixings & Index name, fixing date, value, and source timestamp \\
\bottomrule
\end{tabularx}
\end{center}

Validation must reject inconsistent pay/receive signs, duplicate
exercise dates, exercise after maturity, missing historical fixings,
unsupported FX quote direction, and schedules that extend beyond the
relevant curves.

\subsection{Market fields}\label{42-market-fields}

\begin{itemize}
\tightlist
\item
  domestic discount curve and any distinct domestic forecast curve;
\item
  foreign discount curve and any distinct foreign forecast curve;
\item
  cross-currency basis incorporated consistently in curves or explicit
  collateralized FX forwards;
\item
  domestic normal swaption-volatility surface;
\item
  foreign normal swaption-volatility surface;
\item
  FX spot and Black volatility term structure;
\item
  correlations \passthrough{\lstinline!rho\_dom\_for!},
  \passthrough{\lstinline!rho\_dom\_fx!}, and
  \passthrough{\lstinline!rho\_for\_fx!}; and
\item
  quote timestamps, market snapshot ID, source IDs, units, and
  stale-data flags.
\end{itemize}

The correlation matrix must be symmetric, have ones on the diagonal, and
be positive semidefinite. Reject it if materially invalid; a small
documented eigenvalue floor may repair rounding noise only, and the
result must disclose that repair.

\subsection{Model and numerical
fields}\label{43-model-and-numerical-fields}

\begin{itemize}
\tightlist
\item
  fixed or calibrated Hull-White mean reversions and volatility
  parameterization for each currency;
\item
  calibration helper set and calibration error type;
\item
  FX volatility interpolation/extrapolation rule;
\item
  correlation source and effective date;
\item
  training-path count, pricing-path count, time steps, seed, RNG type,
  antithetic flag, and optional Sobol settings;
\item
  regression basis ID, degree, regularization, standardization rule, and
  candidate-path rule; and
\item
  requested risk scenarios and bump sizes.
\end{itemize}

\section{Market construction and
calibration}\label{5-market-construction-and-calibration}

\subsection{Curves}\label{51-curves}

Build each currency\textquotesingle s discount and forecast curves from
instruments appropriate to that currency and collateral agreement. The
current \passthrough{\lstinline!build\_sofr\_curve()!} is the domestic
USD template: it uses an overnight deposit at the front and SOFR OIS
helpers beyond it. Generalize the pattern only when implementation
begins; do not force non-USD markets through
\passthrough{\lstinline!ql.Sofr!}.

Cross-currency valuation must be internally arbitrage-consistent. In
particular, the domestic and foreign discount curves, spot FX,
collateral convention, cross-currency basis, and model-implied FX
forwards must agree within tolerance. A deterministic non-callable
cross-currency swap is the mandatory first calibration check.

The first model has one stochastic rate factor per currency and no
stochastic basis. The market snapshot must therefore supply:

\begin{itemize}
\tightlist
\item
  the domestic collateral discount curve
  \passthrough{\lstinline!P\_d(0,T)!};
\item
  a foreign effective curve \passthrough{\lstinline!P\_f\_star(0,T)!}
  consistent with domestic collateral, cross-currency basis, spot FX,
  and collateralized FX forwards; and
\item
  any distinct index forecast curve as a deterministic basis mapping to
  the corresponding modeled discount curve.
\end{itemize}

For a forecast curve \passthrough{\lstinline!P\_c\_fwd!}, freeze the
initial multiplicative basis ratio in version 1:

\begin{align*}
  B_c(0,T) &= \frac{P_c^{\mathrm{fwd}}(0,T)}{P_c^{\mathrm{disc}}(0,T)}, \\
  P_c^{\mathrm{fwd}}(t,T)
    &= P_c^{\mathrm{disc}}(t,T)\frac{B_c(0,T)}{B_c(0,t)}.
\end{align*}

Floating coupons derive from \passthrough{\lstinline!P\_c\_fwd(t,T)!}.
This makes the multi-curve assumption explicit: rate uncertainty comes
from the currency\textquotesingle s Hull-White factor, while
forecast/discount and cross-currency basis are deterministic. A
stochastic-basis model is a different model version.

\subsection{Hull-White calibration}\label{52-hull-white-calibration}

Calibrate a separate \passthrough{\lstinline!ql.HullWhite!} model to
each currency\textquotesingle s normal swaption surface. Stay close to
the existing repository pattern:

\begin{itemize}
\tightlist
\item
  create \passthrough{\lstinline!ql.SwaptionHelper!} objects with
  normal-volatility quotes;
\item
  use \passthrough{\lstinline!ql.JamshidianSwaptionEngine!} for helper
  valuation;
\item
  use \passthrough{\lstinline!ql.LevenbergMarquardt!} and explicit end
  criteria;
\item
  fix or calibrate mean reversion according to the model configuration;
  and
\item
  report market value, model value, and calibration error for every
  helper.
\end{itemize}

For the first benchmark, use a diagonal strip aligned to the
trade\textquotesingle s exercise dates and remaining maturity, as
\passthrough{\lstinline!bermudan\_diagonal\_calibration\_pillars()!}
does today. Production validation should also compare against a broader
surface calibration so a diagonal-only fit does not hide off-strip
instability.

\subsection{FX calibration}\label{53-fx-calibration}

Represent FX volatility with QuantLib
\passthrough{\lstinline!BlackVolTermStructure!} objects, starting with
piecewise-constant ATM Black volatilities. "Black-Scholes FX" here means
a lognormal FX diffusion. Because interest rates are stochastic, it does
not mean that the complete product has the deterministic-rate analytic
Black-Scholes price.

Use market Black implied volatilities as initial targets, not
automatically as the hybrid diffusion coefficients. Calibrate the
piecewise \passthrough{\lstinline!sigma\_fx(t)!} buckets by pricing the
same FX vanilla instruments in the two-rate-plus-FX hybrid model, then
solve until model prices or re-implied Black volatilities meet section
11.2. Use the same spot, effective curves, delta, premium, and quote
conventions as the trading system.

\subsection{Correlations}\label{54-correlations}

Correlations are model inputs, not calibration leftovers. Source them
from a governed historical/implied process, record observation windows
and transformations, and provide stress values. At minimum report price
sensitivity to each correlation and to a joint correlation stress.

\section{Joint dynamics}\label{6-joint-dynamics}

Work under the domestic risk-neutral measure. Let
\passthrough{\lstinline!r\_d(t)!} and \passthrough{\lstinline!r\_f(t)!}
be the domestic and foreign short rates and let
\passthrough{\lstinline!X(t) = log(S(t))!}.

Each rate follows a Hull-White 1F process fitted exactly to its modeled
effective curve. Let \passthrough{\lstinline!theta\_f\_for(t)!} be the
foreign-measure drift fitted to
\passthrough{\lstinline!P\_f\_star(0,T)!}. With
\passthrough{\lstinline!S = DOM/FOR!}, the domestic-measure form is:

\begin{align*}
  \mathrm{d}r_d(t)
    &= \left[\theta_d(t)-a_dr_d(t)\right]\mathrm{d}t
       +\sigma_d(t)\mathrm{d}W_d(t), \\
  \mathrm{d}r_f(t)
    &= \left[\theta_f^{\mathrm{for}}(t)-a_fr_f(t)\right]\mathrm{d}t \\
    &\quad-\rho_{f,\mathrm{FX}}\sigma_f(t)\sigma_{\mathrm{FX}}(t)\mathrm{d}t
       +\sigma_f(t)\mathrm{d}W_f(t).
\end{align*}

FX follows a Black-Scholes lognormal process under the domestic measure:

\begin{equation*}
  \frac{\mathrm{d}S_t}{S_t}
    = \left[r_d(t)-r_f(t)\right]\mathrm{d}t
      +\sigma_{\mathrm{FX}}(t)\mathrm{d}W_{\mathrm{FX}}(t).
\end{equation*}

The Brownian increments satisfy the configured 3x3 correlation matrix.
The displayed quanto term assumes the stated FX quotation, domestic
money-market numeraire, scalar factor loadings, and the sign convention
above. The implementation must derive and unit-test the drift from those
conventions; changing the FX orientation or numeraire changes the
formula. It is not valid to simulate an unadjusted foreign-measure
Hull-White process independently and then attach FX afterward.

The implementation must document the exact numeraire and FX quote
convention in code and model metadata. Martingale tests must verify
discounted domestic assets and the domestic value of foreign
money-market/zero-coupon assets within Monte Carlo error.

\section{Simulation design}\label{7-simulation-design}

\subsection{Time grid}\label{71-time-grid}

The grid must contain every fixing, accrual boundary needed by the
payoff, payment, notional exchange, notice, decision, effective
exercise, and settlement date. Add intermediate steps so no interval
exceeds the configured maximum, initially one month. OIS coupons should
use QuantLib\textquotesingle s compounded-overnight cash-flow
conventions or a validated conditional approximation; do not simulate
every overnight fixing unless the accuracy test proves it necessary.

\subsection{Path generation}\label{72-path-generation}

\begin{itemize}
\tightlist
\item
  use exact conditional Gaussian transitions for the Hull-White factors
  where available;
\item
  use a log update for FX so simulated FX stays positive;
\item
  correlate independent normal draws with a Cholesky or eigen
  factorization of the validated correlation matrix;
\item
  integrate short rates consistently for pathwise domestic discount
  factors;
\item
  use antithetic paths initially and optionally Brownian bridge plus
  Sobol after pseudo-random validation; and
\item
  split training and pricing samples so the same paths do not both fit
  and value the stopping rule.
\end{itemize}

QuantLib 1.41 in the current Python environment exposes
\passthrough{\lstinline!ql.HullWhiteProcess!},
\passthrough{\lstinline!ql.BlackScholesMertonProcess!},
\passthrough{\lstinline!ql.StochasticProcessArray!}, and
\passthrough{\lstinline!ql.GaussianMultiPathGenerator!}. Do not simply
place two independent rate processes and
\passthrough{\lstinline!ql.BlackScholesMertonProcess!} in an array: that
does not automatically create the pathwise
\passthrough{\lstinline!r\_d-r\_f!} FX drift or the foreign-rate quanto
adjustment. Use an explicit joint evolver (or custom QuantLib process)
whose transition equations implement section 6;
QuantLib\textquotesingle s grids, random sequences, correlations, curve
objects, and statistics remain reusable.

\subsection{Pathwise cash-flow
value}\label{73-pathwise-cash-flow-value}

On each path:

\begin{enumerate}
\def\labelenumi{\arabic{enumi}.}
\tightlist
\item
  generate both rate factors, log FX, and domestic discount factors;
\item
  determine all already-fixed and simulated floating coupons without
  look-ahead;
\item
  convert each foreign payment to domestic currency at the contractually
  applicable FX rate on that path;
\item
  include or exclude notional exchanges exactly as the trade specifies;
\item
  apply call event ordering; and
\item
  discount realized domestic-equivalent cash flows to the valuation
  date.
\end{enumerate}

At an exercise date, the state-dependent value of remaining legs should
use analytic Hull-White conditional zero-coupon bond prices for
\passthrough{\lstinline!P\_d!} and \passthrough{\lstinline!P\_f\_star!}
where possible. Distinct forecast bonds come from the deterministic
basis mapping in section 5.1. This reduces inner simulation and keeps
the method a single outer Monte Carlo with regression, not a slow nested
Monte Carlo.

\section{Early exercise decision}\label{8-early-exercise-decision}

\subsection{Decision rule}\label{81-decision-rule}

For each alive path \passthrough{\lstinline!p!} and exercise date
\passthrough{\lstinline!t\_i!}, define exclusive decision components as
in section 2.5:

\begin{itemize}
\tightlist
\item
  \passthrough{\lstinline!E\_i(p)!} - exercise-only value
  (\passthrough{\lstinline!J\_i!} for cancellation, or entered-swap
  value plus settlement for entry); and
\item
  \passthrough{\lstinline!C\_i(p)!} - the conditional expected value of
  the mutually exclusive continuation cash flows and later optionality,
  estimated by regression.
\end{itemize}

Exercise when:

\begin{equation*}
  E_i(p)>C_i(p)+\varepsilon_{\mathrm{ex}}.
\end{equation*}

Use a deterministic tie rule: continue when the values are equal within
tolerance. Add the common surviving component
\passthrough{\lstinline!A\_i!} after the decision comparison. For a
zero-fee cancellation, \passthrough{\lstinline!E\_i = J\_i = 0!}; the
holder cancels when the estimated value of cancellable continuation cash
flows and later optionality is negative enough.

\subsection{Backward Longstaff-Schwartz
algorithm}\label{82-backward-longstaff-schwartz-algorithm}

\begin{enumerate}
\def\labelenumi{\arabic{enumi}.}
\tightlist
\item
  Simulate training paths and compute terminal values.
\item
  Starting from the last call date, regress the discounted realized
  value from continuing on state variables observed at that date.
\item
  Compare \passthrough{\lstinline!E\_i!} with the fitted
  \passthrough{\lstinline!C\_i!} and replace the future path value with
  the exercise value when the decision rule triggers.
\item
  Move one call date backward and repeat using only information
  available at that date.
\item
  Freeze all preprocessing, coefficients, and exercise rules.
\item
  Simulate an independent pricing sample, apply the frozen policy, and
  average discounted stopped cash flows.
\end{enumerate}

This is the same economic comparison performed by a tree during backward
induction, but the conditional continuation value is estimated from
simulated cross-currency states rather than read from lattice nodes.

\subsection{Regression state and
basis}\label{83-regression-state-and-basis}

Initial standardized state variables:

\begin{itemize}
\tightlist
\item
  domestic Hull-White factor or short rate;
\item
  foreign Hull-White factor or short rate;
\item
  \passthrough{\lstinline!log(S)!};
\item
  domestic value of the remaining non-callable swap;
\item
  remaining domestic-leg PV;
\item
  domestic value of remaining foreign-leg PV; and
\item
  time to final maturity.
\end{itemize}

Initial basis: constant, linear terms, squares, and pairwise
cross-products. Prefer orthogonal polynomials after standardization and
cap the basis at a predeclared size. Record rank, condition number,
coefficient vector, training count, and out-of-sample residual
diagnostics at every exercise date. If regression is ill-conditioned or
under-populated, fail with a diagnostic or use a preapproved simpler
basis; never silently return an exercise-always rule.

Use all alive paths for a cancellation right unless validation supports
a narrower candidate set. The ordinary "in-the-money paths only"
shortcut is not directly applicable when the immediate cancellation
payoff is zero and the economic benefit is avoiding a negative
continuation value.

\section{QuantLib integration
plan}\label{9-quantlib-integration-plan}

\subsection{Reuse from this
repository}\label{91-reuse-from-this-repository}

Preserve and reuse these design patterns when implementation is
authorized:

\begin{itemize}
\tightlist
\item
  \passthrough{\lstinline!valuation\_date()!} and the single
  evaluation-date discipline;
\item
  explicit calendars, business-day handling, schedules, day counts, and
  spot starts;
\item
  OIS construction and calibration repricing checks from
  \passthrough{\lstinline!build\_sofr\_curve()!};
\item
  \passthrough{\lstinline!ql.SwaptionHelper!} construction and
  normal-vol conversions;
\item
  Hull-White calibration and parameter reporting from
  \passthrough{\lstinline!build\_bermudan\_short\_rate\_model()!};
\item
  trade-state normalization, structured result rows, deterministic
  defaults, and bump-and-revalue reporting; and
\item
  unit-test style in
  \passthrough{\lstinline!tests/test\_sofr\_curve.py!}.
\end{itemize}

\subsection{Do not reuse as the cross-currency
engine}\label{92-do-not-reuse-as-the-cross-currency-engine}

\begin{itemize}
\tightlist
\item
  \passthrough{\lstinline!ql.Swaption!} represents a single underlying
  swaption, not this two-currency callable cash-flow set.
\item
  \passthrough{\lstinline!ql.TreeSwaptionEngine!} is a short-rate
  lattice engine for swaptions. It does not evolve two currency curves
  plus FX.
\item
  a single domestic Hull-White tree with deterministic FX forwards omits
  foreign-rate and FX/rate correlation risks and is only an
  approximation/control, not the specified model.
\end{itemize}

\subsection{Python reference versus C++
production}\label{93-python-reference-versus-c-production}

QuantLib\textquotesingle s C++ library provides generic Monte Carlo and
Longstaff-Schwartz infrastructure, but the installed Python 1.41 wrapper
does not expose a generic
\passthrough{\lstinline!LongstaffSchwartzPathPricer!}, and neither the
installed package nor the current repo has a callable cross-currency
instrument. Therefore:

\begin{itemize}
\tightlist
\item
  build the first transparent reference pricer in this repository with
  QuantLib market/process objects and an explicit, testable LSM layer
  only after approval; and
\item
  implement the production path pricer/engine in C++ against a pinned
  QuantLib version, using QuantLib cash flows, curves, processes, random
  sequences, regression utilities, and pricing-engine conventions.
\end{itemize}

Do not make a development branch of QuantLib itself the production
dependency. The current Python environment is QuantLib 1.41. QuantLib
1.43 is the latest official release as of this specification and adds
constant-notional cross-currency swap instruments, which should be
evaluated for the deterministic C++ layer. It still does not supply the
complete callable three-factor MC-LSM engine specified here. Pin the
selected release only after compatibility and regression testing.

\section{Outputs}\label{10-outputs}

Every result must include:

\begin{itemize}
\tightlist
\item
  NPV in reporting currency and currency unit;
\item
  standard error and 95 percent confidence interval;
\item
  non-callable underlying NPV;
\item
  callable value and embedded option value under a documented sign
  convention;
\item
  path counts, rejected paths, time-grid size, seed, RNG, antithetic
  setting, and runtime;
\item
  model parameters, calibration rows, correlation matrix, and market
  snapshot ID;
\item
  unconditional exercise probability by date, conditional exercise
  probability among surviving paths, no-exercise probability, and
  expected exercise date;
\item
  regression diagnostics by date;
\item
  cash-flow PV by leg and currency; and
\item
  requested curve, volatility, FX spot, correlation, and model-parameter
  risks.
\end{itemize}

All values are from the unilateral exercise holder\textquotesingle s
perspective. Define
\passthrough{\lstinline!embedded\_option\_value = callable\_value - non\_callable\_value!};
it must be nonnegative within the comparison error defined in section 11
because the holder can always choose not to exercise. For the
counterparty perspective, negate the full callable trade result or price
a separately defined request with the opposite holder; do not reverse
cash-flow signs while silently preserving option ownership.

\section{Validation and acceptance
criteria}\label{11-validation-and-acceptance-criteria}

For every comparison below, define
\passthrough{\lstinline!comparison\_standard\_error!} explicitly:

\begin{itemize}
\tightlist
\item
  on common-random-number or otherwise paired paths, compute the
  pathwise payoff difference
  \passthrough{\lstinline!D\_p = PV\_A(p) - PV\_B(p)!} and use
  \passthrough{\lstinline!stdev(D\_p) / sqrt(N)!};
\item
  on independent samples, use
  \passthrough{\lstinline!sqrt(SE\_A\^2 + SE\_B\^2)!}; and
\item
  if a more general correlated estimator is used, use
  \passthrough{\lstinline!sqrt(SE\_A\^2 + SE\_B\^2 - 2 * covariance(mean\_A, mean\_B))!}
  and report the covariance estimate.
\end{itemize}

Do not add marginal standard errors or ignore covariance in a paired
comparison.

\subsection{Deterministic and limiting
tests}\label{111-deterministic-and-limiting-tests}

\begin{itemize}
\tightlist
\item
  With no call dates, equal the analytic non-callable XCCY PV within
  \passthrough{\lstinline!max(3 * comparison\_standard\_error, 0.25 bp * domestic\_notional)!}.
\item
  With FX volatility and both rate volatilities near zero, match
  deterministic discounted cash flows within
  \passthrough{\lstinline!0.10 bp * domestic\_notional!}.
\item
  With identical currencies, unit FX, identical curves, and consistent
  legs, reduce to the corresponding single-currency result within
  \passthrough{\lstinline!max(3 * comparison\_standard\_error, 0.25 bp * domestic\_notional)!}.
\item
  With one exercise date, match a high-accuracy European or nested-MC
  benchmark within
  \passthrough{\lstinline!max(3 * comparison\_standard\_error, 0.50 bp * domestic\_notional)!}.
\item
  With a prohibitively expensive exercise fee, converge to the
  non-callable value.
\item
  With an immediately dominant zero-fee cancellation, the first-date
  conditional exercise probability must be at least 99 percent.
\item
  \passthrough{\lstinline!callable\_value - non\_callable\_value!} must
  equal the reported embedded option value and be no less than
  \passthrough{\lstinline!-3 * comparison\_standard\_error!}.
\end{itemize}

\subsection{Martingale and calibration
tests}\label{112-martingale-and-calibration-tests}

\begin{itemize}
\tightlist
\item
  Reprice both OIS calibration sets to no more than
  \passthrough{\lstinline!0.01 bp!} par-rate error and
  \passthrough{\lstinline!1e-8 * notional!} absolute NPV error.
\item
  Reprice domestic and foreign swaption helpers to no more than 10
  percent relative price error per helper and 5 percent equal-weight RMS
  relative price error across the approved strip. For helper
  \passthrough{\lstinline!i!},
  \passthrough{\lstinline!e\_i = V\_model\_i / V\_market\_i - 1!}, and
  the implemented strip metric is
  \passthrough{\lstinline!sqrt(sum(e\_i\^2) / N)!}.
\item
  Reprice FX vanilla calibration instruments to no more than
  \passthrough{\lstinline!0.25!} volatility percentage points of
  re-implied Black-vol error per quote.
\item
  Validate domestic-measure martingales for domestic and converted
  foreign zero-coupon assets within three standard errors and
  \passthrough{\lstinline!0.25 bp!} of notional.
\item
  Reprice a deterministic cross-currency swap against an independent
  cash-flow calculation within
  \passthrough{\lstinline!0.10 bp * domestic\_notional!}.
\end{itemize}

\subsection{Numerical convergence}\label{113-numerical-convergence}

Demonstrate stability across:

\begin{itemize}
\tightlist
\item
  at least three path counts;
\item
  at least two time-grid densities;
\item
  multiple independent seeds;
\item
  simpler and richer regression bases;
\item
  pseudo-random and, after validation, low-discrepancy sequences; and
\item
  training/pricing sample splits.
\end{itemize}

For adjacent path-count or grid refinements, require an absolute price
difference no greater than
\passthrough{\lstinline!max(3 * comparison\_standard\_error, 0.50 bp * domestic\_notional)!}.
For the official pricing sample, require:

\begin{equation*}
  \mathrm{SE}
  \leq
  \min\!\left(
    1.00\,\mathrm{bp}\,N_d,
    \max\!\left(0.5\%\,\lvert V_{\mathrm{option}}\rvert,
                 0.05\,\mathrm{bp}\,N_d\right)
  \right).
\end{equation*}

Pricing-path error is conditional on the fitted exercise policy and does
not include training-policy uncertainty or LSM lower-bound bias. Run at
least five independent training seeds. Freeze one common pricing sample,
with at least the official pricing-path count and satisfying the
single-price error threshold above, and evaluate every trained policy on
those identical paths. The standard deviation across these paired policy
values estimates training-policy dispersion and must be no greater than
\passthrough{\lstinline!0.50 bp * domestic\_notional!}; report it
separately from the pricing-sample standard error. Do not use
independently seeded pricing samples for this test because that would
mix ordinary pricing noise with policy uncertainty.

\subsection{Independent benchmark}\label{114-independent-benchmark}

Model validation should compare at least one trade with an independent
implementation or vendor. A simplified two- or three-dimensional
PDE/lattice benchmark is useful for short maturities and sparse calls
even though it is not the production engine.

\section{Risks and controls}\label{12-risks-and-controls}

\begin{center}
\begin{tabularx}{\columnwidth}{>{\bfseries}p{0.38\columnwidth}X}
\toprule
Risk & Required control \\
\midrule
Wrong FX quotation or sign & Explicit \code{DOM/FOR} metadata and unit
  tests for conversion \\
Wrong measure or quanto drift & Martingale tests and documented numeraire \\
Regression look-ahead bias & Independent training and pricing paths \\
Overfitting & Fixed basis policy, conditioning diagnostics, and
  out-of-sample tests \\
Hidden curve inconsistency & Deterministic XCCY repricing and FX-forward
  checks \\
Correlation instability & PSD validation, source metadata, and
  bump/stress results \\
Date/event ambiguity & Four-date exercise records and explicit event ordering \\
Global QuantLib evaluation date & One valuation context per job/process
  and concurrency tests \\
False numerical precision & Standard error, convergence report, and
  rounded display \\
\bottomrule
\end{tabularx}
\end{center}

\section{Proposed implementation sequence - future work
only}\label{13-proposed-implementation-sequence---future-work-only}

\begin{enumerate}
\def\labelenumi{\arabic{enumi}.}
\tightlist
\item
  Freeze the trade and market schemas and build a deterministic
  non-callable XCCY cash-flow pricer.
\item
  Add both calibrated Hull-White models and FX vanilla validation.
\item
  Add the joint domestic-measure simulator and martingale tests.
\item
  Add one-date exercise and an independent European benchmark.
\item
  Add two-pass Bermudan LSM, exercise diagnostics, and convergence
  tests.
\item
  Add common-random-number bump risks and batch performance work.
\item
  Port the validated engine to the C++ integration boundary in the
  companion guide.
\end{enumerate}

No phase may start by changing the existing single-currency
\passthrough{\lstinline!TreeSwaptionEngine!} path.

\section{References}\label{14-references}

\begin{itemize}
\tightlist
\item
  QuantLib official documentation:
  \url{https://www.quantlib.org/docs.shtml}
\item
  QuantLib releases, including 1.43:
  \url{https://github.com/lballabio/QuantLib/releases}
\item
  QuantLib 1.43 \passthrough{\lstinline!TreeSwaptionEngine!}:
  \url{https://github.com/lballabio/QuantLib/blob/v1.43/ql/pricingengines/swaption/treeswaptionengine.hpp}
\item
  QuantLib 1.43 exercise types:
  \url{https://github.com/lballabio/QuantLib/blob/v1.43/ql/exercise.hpp}
\item
  QuantLib 1.43 Longstaff-Schwartz path pricer:
  \url{https://github.com/lballabio/QuantLib/blob/v1.43/ql/methods/montecarlo/longstaffschwartzpathpricer.hpp}
\item
  QuantLib 1.43 Monte Carlo base:
  \url{https://github.com/lballabio/QuantLib/blob/v1.43/ql/pricingengines/mcsimulation.hpp}
\item
  Longstaff, F. A. and Schwartz, E. S. (2001), "Valuing American Options
  by Simulation: A Simple Least-Squares Approach," Review of Financial
  Studies 14(1), 113-147.
\end{itemize}

\section{Rollback}\label{rollback}

This is a specification only. Removing this LaTeX source and its
generated PDF fully rolls back the document change; application behavior
and pricing results are unaffected.

\end{document}
